\documentclass[10pt]{article}
\usepackage[T1]{fontenc}
\usepackage[utf8]{inputenc}
\usepackage[french]{babel}
\usepackage{amsmath,amssymb,mathrsfs,mathpazo}
\usepackage{geometry}
\geometry{a4paper,margin=2cm}
\usepackage{graphicx,array,tabularx,tabularray,colortbl,multirow,multicol,booktabs,enumitem,pifont,ulem,fancybox,diagbox,cancel}
\usepackage{pstricks,pstricks-add,pst-plot,pst-tree,pst-node,pst-text,pst-func,pst-math,pst-bspline,pst-3dplot,pst-eucl,pst-solides3d}
\let\pscalculate\undefined
\usepackage{pst-bezier}
\usepackage{tikz,pgfplots,tkz-tab}
\usetikzlibrary{arrows,arrows.meta,patterns,calc,positioning,trees,decorations.markings,shapes.geometric,shapes.misc,shapes.arrows,intersections,decorations.pathmorphing,decorations.pathreplacing,backgrounds,shadows,babel}
\pgfplotsset{compat=1.18}
\usepgfplotslibrary{fillbetween}
\usepackage[np]{numprint}
\usepackage{listings,makecell,pas-tableur,twoopt}
\let\barre\undefined
\usepackage[pstricks]{tablvar}
\definecolor{darkgreen}{RGB}{0,100,0}
\providecommand{\red}{\color{red}}
\providecommand{\blue}{\color{blue}}
\providecommand{\green}{\color{green}}
\providecommand{\black}{\color{black}}
\providecommand{\white}{\color{white}}
\providecommand{\gray}{\color{gray}}
\providecommand{\cyan}{\color{cyan}}
\providecommand{\magenta}{\color{magenta}}
\providecommand{\yellow}{\color{yellow}}
\providecommand{\violet}{\color{violet}}
\providecommand{\orange}{\color{orange}}
\providecommand{\euro}{\texteuro}
\providecommand{\EUR}{\texteuro}
\providecommand{\ieme}{\textsuperscript{e}}
\providecommand{\degres}{\ensuremath{^\circ}}
\providecommand{\degree}{\ensuremath{^\circ}}
\providecommand{\vv}[1]{\overrightarrow{#1}}
\providecommand{\Coord}[1]{#1}
\providecommand{\decosix}{}
\providecommand{\starredbullet}{\ensuremath{\bullet}}
\pagestyle{empty}
\setlength{\parindent}{0pt}
\usepackage{adjustbox}\DeclareUnicodeCharacter{2236}{:}\DeclareUnicodeCharacter{2190}{\ensuremath{\leftarrow}}
\providecommand{\R}{\mathbb{R}}
\providecommand{\N}{\mathbb{N}}
\providecommand{\Z}{\mathbb{Z}}
\providecommand{\Q}{\mathbb{Q}}
\providecommand{\D}{\mathbb{D}}
\providecommand{\C}{\mathbb{C}}
\providecommand{\e}{{\mathrm{e}}}
\providecommand{\vect}[1]{\overrightarrow{#1}}
\providecommand{\vectt}[1]{\overrightarrow{\mathrm{#1}}}
\providecommand{\vvt}[1]{\overrightarrow{\mathrm{#1}}}
\providecommand{\np}[2][]{{#2~#1}}
\providecommand{\vv}[1]{{\overrightarrow{#1}}}
\renewcommand{\d}{{\mathrm{d}}}
\providecommand{\degres}{{^\circ}}
\providecommand{\degre}{{^\circ}}
\providecommand{\blue}{\color{blue}}
\providecommand{\red}{\color{red}}
\providecommand{\enskip}{\quad}
\providecommand{\llbracket}{\lbrack\!\lbrack}
\providecommand{\rrbracket}{\rbrack\!\rbrack}
\providecommand{\euro}{{\text{€}}}
\providecommand{\no}{{\text{n°}}}
\providecommand{\interval}[2]{{[#1;#2]}}
\providecommand{\intervFO}[2]{{[#1;#2[}}
\providecommand{\dotfill}{\ldots}

\providecommand{\ds}{\displaystyle}
\providecommand{\textding}[1]{#1}
\providecommand{\Oij}{\(\mathrm{O};\vec{\imath},\vec{\jmath}\)}
\providecommand{\Oijk}{\(\mathrm{O};\vec{\imath},\vec{\jmath},\vec{k}\)}
\providecommand{\Ouv}{\(\mathrm{O};\vec{u},\vec{v}\)}
\providecommand{\intervFF}[2]{\ensuremath{[#1;#2]}}
\renewcommand{\euro}{\ifmmode\text{\texteuro}\else\texteuro\fi}
\DeclareUnicodeCharacter{00B0}{\ifmmode^{\circ}\else\textdegree\fi}
\DeclareUnicodeCharacter{20AC}{\ifmmode\text{\texteuro}\else\texteuro\fi}
% Énoncés : Éducation nationale. Transcription et corrigés : APMEP, Denis Vergès et contributeurs. Découpage et classement : CoopMaths, Sébastien Lozano, Éric Elter et contributeurs. Adaptation de la présentation : Maths Explorer.
% CC BY-SA : https://coopmaths.fr/www/about/
% https://coopmaths.fr/alea/static/bac/2025/tex/bac_2025_06_sujet1_etrangers_4_cor.tex
\begin{document}
\noindent\textit{APMEP / CoopMaths. CC BY-SA. Présentation adaptée par Maths Explorer.}\par\medskip
\medskip

\textbf{Partie A : }
\medskip

On considère l'équation différentielle $\left(E_{1}\right): y'+0,48 y=\dfrac{1}{250}$ où $y$ est une fonction de la variable $t$ définie sur l'intervalle $[0 ;+\infty[$.

\begin{enumerate}
\item La fonction $h$ est dérivable  sur $[0 ;+\infty[$ et pour  tout réel positif $t$, $h'(t)=0$ 
On a dont $h'(t) +0,48 h(t) =0 + 0,48 \times \dfrac{1}{120}= \dfrac{1}{250}$ 

Donc  la fonction $h$ est solution de l'équation différentielle $\left(E_{1}\right)$.
\item Les solutions de l'équation différentielle $y'+0,48 y=0$ sont les fonctions de la forme $f(t) =C \e^{-0,48t}$ avec $C \in \R$. 
\item Les solutions de l'équation différentielle $\left(E_{1}\right)$ sont donc les fonctions de la forme $f(t) =C \e^{-0,48t} + \dfrac{1}{120}$ avec $C \in \R$.
\end{enumerate}
\medskip

\textbf{Partie B :}
\medskip

\begin{enumerate}
\item Supposons que $p$ est solution de l'équation différentielle ( $E_{2}$ ),

On a donc $p'=\dfrac{1}{250} p \times(120-p)$ 

Or $p=\dfrac{1}{y}$ et $p'= \dfrac{-y'}{y^2}$ 

D'où  $\begin{aligned} \dfrac{-y'}{y^2}= \dfrac{1}{250} \dfrac{1}{y} \times\left( 120-\dfrac{1}{y}\right)& \Longrightarrow -y'=y^2 \times  \dfrac{1}{250y }\times \dfrac{120y-1}{y}\\
&\Longleftrightarrow -y' = \dfrac{1}{250} \times (120y-1)\\
&\Longleftrightarrow -y'= 0,48 y -\dfrac{1}{250}\\
&\Longleftrightarrow y'+0,48 y=\dfrac{1}{250}
\end{aligned}$.

Donc $y$ est solution de l'équation différentielle $\left(E_{1}\right): y'+0,48 y=\dfrac{1}{250}$.

\item $y$ est de la forme $y(t)=C \e^{-0,48t} + \dfrac{1}{120}$ avec $C \in \R$.

on a donc $\begin{aligned}[t] p(t)& = \dfrac{1}{y(t)}\\
& = \dfrac{1}{C \e^{-0,48t} + \frac{1}{120}} \, \text{ avec } C \in \R\\
& = \dfrac{120}{120 C \e^{-0,48t} +1 } \, \text{ avec } C \in \R\\
& = \dfrac{120}{K \e^{-0,48t} +1 } \, \text{ avec } K \in \R \text { en posant }  K=120C\\
\end{aligned}$

\item  On sait que  $p(0)=30$ donc $\dfrac{120}{K \e^{0} +1 }=30$ 

donc $\dfrac{120}{30}= K+1$   d'où $K =\dfrac{120}{30}-1=3$

\item On sait que : $\begin{cases}
\lim\limits_{t \to +\infty }-0,48t= -\infty \\
\lim\limits_{X\to -\infty} \e^X = 0
\end{cases}$ donc, par composition, $\lim\limits_{t\to +\infty }\e^{-0,48t}=0$ 

Donc, par produit et somme, $\lim\limits_{t \to +\infty}1+3\e^{-0,48t}=1$ 

et finalement par quotient, $\lim\limits_{t \to +\infty}p(t) =120$ 

Dans le contexte de l'exercice la population de la bactérie  se stabilisera vers $\np{120000}$. 



\item $\begin{aligned}[t] p(t)=60 & \iff \dfrac{120}{3 \e^{-0,48t} +1 }=60 \\
	& \iff 120 = 60 \times \left(3 \e^{-0,48t} +1\right)\\
	& \iff 2 = 3 \e^{-0,48t} +1\\
	& \iff \dfrac{1}{3}= \e^{-0,48t}\\
	& \iff -0,48 t = \ln\left(\dfrac{1}{3}\right)=-\ln(3)\\
	& \iff t = \dfrac{\ln(3)}{0,48}\approx 2,289
	\end{aligned}$
	
Le temps nécessaire pour que la population de bactéries dépasse $\np{60000}$ individus est de $2,289$ heures soit 2 heures et 17 minutes.
\end{enumerate}

\end{document}
